\documentclass[11pt]{article}

\usepackage[margin=1in]{geometry}
\usepackage[T1]{fontenc}
\usepackage[utf8]{inputenc}
\usepackage{lmodern}
\usepackage{amsmath,amssymb}
\usepackage{booktabs}
\usepackage{array}
\usepackage{enumitem}
\usepackage{longtable}
\usepackage{xcolor}
\usepackage{hyperref}
\usepackage{listings}

\hypersetup{
    colorlinks=true,
    linkcolor=blue,
    urlcolor=blue
}

\lstset{
    basicstyle=\ttfamily\small,
    breaklines=true,
    columns=fullflexible,
    keepspaces=true,
    frame=single
}

\title{Project Documentation Index for \texttt{dev}}
\author{}
\date{}

\begin{document}
\maketitle

\section{Purpose of this Folder}

This folder is a compact but human-readable description of the current \texttt{dev} codebase. It is meant to help a reader, or an LLM that is about to modify the repository, understand the project at the level of design intent, invariants, execution flow, and branch-specific semantics.

It is \emph{not} a full thesis-style documentation set. It deliberately focuses on the parts that are easy to miss when reading source code line by line:
\begin{itemize}[leftmargin=2em]
    \item what the project is trying to compare,
    \item why the code is organized the way it is,
    \item what assumptions the experiment runner enforces,
    \item which branch semantics are special,
    \item and which project contracts should be preserved when updating the code.
\end{itemize}

\section{Important Reading Notes}

This folder describes the \textbf{latest current version} of the project only. It does not try to preserve older superseded behavior.

A few global facts are worth knowing before reading the rest:
\begin{itemize}[leftmargin=2em]
    \item The current top-level entry point is \texttt{python -m dev.main}.
    \item The active experiment driver is the generalized study runner under \texttt{dev/experiments/study/}.
    \item \texttt{dev/master\_config.py} is the single intended user-editable configuration file.
    \item Dynamic branches prepare one shared dynamic map state for the whole branch execution.
    \item The experiment compares \textbf{classical mapping} and \textbf{cyclic mapping}, not two different planning algorithms.
\end{itemize}

\section{Recommended Reading Order}

\begin{enumerate}[leftmargin=2em]
    \item \texttt{01\_overview} for the big picture.
    \item \texttt{02\_repository\_map} for where the active code lives.
    \item \texttt{03\_composite\_grid} and \texttt{04\_mapping\_system} for the core representation.
    \item \texttt{05\_solver\_behavior} and \texttt{06\_experiment\_protocol} for execution semantics.
    \item \texttt{07\_static\_artificial\_branch}, \texttt{08\_dynamic\_shared\_pipeline}, \texttt{09\_dynamic\_port\_branch}, and \texttt{10\_dynamic\_campus\_branches} for branch-specific details.
    \item \texttt{11\_outputs\_and\_visualization} and \texttt{12\_current\_configuration} for current output and config contracts.
    \item \texttt{13\_update\_guidelines} and \texttt{14\_troubleshooting} before making changes.
\end{enumerate}

\section{Scope Boundary}

These documents assume the reader can still inspect the code. Therefore, they do not duplicate every implementation detail. The aim is to expose the \emph{shape} of the system, the \emph{meaning} of its branches and metrics, and the \emph{assumptions that should survive future updates}.

\end{document}
